\chapter{Mathematical Formulation and Graph Model}
\label{chap:model}
\label{chap:formulation}

The mathematical formulation connects timetable selection, vehicle chaining,
fleet assignment, battery feasibility, charging, parking capacity, and cost
calculation. Solving all of these decisions in one exact model would be very
large. The formulation is therefore used to derive the graph and matching
subproblems solved inside the matheuristic. Here \(\pi\) denotes the selected
timetable, \(G_\pi\) the compatibility graph, \(B\) the set of vehicle blocks,
and \(C^{VS}\) the reported vehicle-scheduling cost.

\section{Notation and Planning Objects}

\subsection{Sets, Parameters, and Indices}

\Cref{tab:notation} summarizes the notation used in the mathematical chapters.

\begin{table}[htbp]
\centering
\caption{Main notation used in the mathematical formulation.}
\label{tab:notation}
\scriptsize
\setlength{\tabcolsep}{3pt}
\renewcommand{\arraystretch}{1.16}
\begin{tabularx}{\textwidth}{@{}>{\RaggedRight\arraybackslash}p{3.75cm}
>{\RaggedRight\arraybackslash}X
>{\RaggedRight\arraybackslash}p{1.05cm}@{}}
\toprule
Symbol & Meaning & Unit \\
\midrule
\(T,D,N\) & Sets of candidate trips, directions, and nodes. & -- \\
\(V,V_E\) & Sets of vehicle types and electric vehicle types, with \(V_E\subseteq V\). & -- \\
\(i,j,b,a,v\) & Indices for trips, vehicle blocks, activities, and vehicle types. & -- \\
\(\pi_d,\pi\) & Selected timetable of direction \(d\), and complete selected timetable. & -- \\
\(st(i),et(i),a(i)\) & Trip start time, trip end time, and main-stop time. & s \\
\(sn(i),en(i)\) & Start and end node of trip \(i\). & -- \\
\(\ell_i\) & Passenger-trip distance. & km \\
\(\ell_a^{dh}\) & Deadhead distance of activity or arc \(a\). & km \\
\(\tau_i^{trip}\) & Passenger-trip driving time \(et(i)-st(i)\). & s \\
\(\tau_a^{dh}\) & Deadhead driving time. & s \\
\(\omega(t)\) & Time-window index containing time \(t\). & -- \\
\(h^{max}_{d,\omega}\) & Maximum headway for direction \(d\) in time window \(\omega\). & s \\
\(\delta^{min}_{n,\omega}\), \(\delta^{max}_{n,\omega}\) & Minimum and maximum stopping time at node \(n\). & s \\
\(t_b^{pi}\), \(t_b^{po}\), \(d(b)\) & Pull-in duration, pull-out duration, and passenger-service duration used for the CO\(_2\) term. & s \\
\(B,G_\pi,H_d\) & Vehicle-block set, compatibility graph, and headway graph. & -- \\
\(u_{bv}\) & Binary vehicle-type variable; equals 1 if block \(b\) uses vehicle type \(v\). & -- \\
\(z_{ij}\) & Binary matching variable; equals 1 if trip \(j\) directly follows trip \(i\). & -- \\
\(y_{ib}\) & Binary coverage variable; equals 1 if trip \(i\) is served by block \(b\). & -- \\
\(r_{ba}^{pre}\), \(r_{ba}^{post}\) & Residual EV autonomy before and after activity \(a\) in block \(b\). & km \\
\(q_a^{slow}\), \(q_a^{fast}\) & Slow and fast charging duration in break activity \(a\). & s \\
\(o_a^{park}\), \(o_a^{slow}\), \(o_a^{fast}\) & Resource-occupancy variables for parking, slow charging, and fast charging. & -- \\
\(g_a\) & Binary or state variable indicating active charging in break window \(a\). & -- \\
\(p_a\) & Paid-break duration of break activity \(a\). & s \\
\(A_v^{max}\), \(T_v^{slow}\), \(q_v^{min,m}\) & EV autonomy, full slow-charge time, and minimum charging duration. & mixed \\
\(\phi_v\) & Fast-charging multiplier for EV type \(v\). & -- \\
\(K_n^{park},K_n^{slow},K_n^{fast}\) & Parking, slow-charge, and fast-charge capacities at node \(n\). & vehicles \\
\(c_v^{fix}\), \(c^{break}\), \(c_v^{pio}\), \(c_v^{\mathrm{CO}_2}\) & Fixed, paid-break, pull-in/out, and CO\(_2\) cost coefficients. & mixed \\
\(C^{VS}\), \(C^{fix}\), \(C^{break}\), \(C^{pull}\), \(C^{\mathrm{CO}_2}\) & Total vehicle-scheduling cost and its fixed, paid-break, pull, and CO\(_2\) components. & cost \\
\(UB\), \(LB^{global}\), \(LB^{TT}(\pi)\) & Validated upper bound, global lower bound, and selected-timetable diagnostic lower bound. & cost \\
\(N_{\mathrm{veh}}\), \(s_{\mathrm{EV}}\), \(\eta_{\mathrm{veh}}\), \(\bar c_T\) & Used vehicles, EV share, vehicle productivity, and average cost per selected trip. & mixed \\
\(M^{veh}\), \(M^{break}\) & Vehicle-first matching constant and paid-break linearization constant. & s \\
\(Q_h(P)\) & Timetable tie-breaking score for path \(P\) under start \(h\). & score \\
\bottomrule
\end{tabularx}
\end{table}
\FloatBarrier

\subsection{Layered Decision Structure}

The solver builds a solution through a sequence of related objects:

\begin{equation}
T
\longrightarrow
\pi
\longrightarrow
G_\pi=(\pi,A_\pi)
\longrightarrow
B
\longrightarrow
(\Omega,C^{VS}).
\label{eq:layered-decision-chain}
\end{equation}

\Cref{eq:layered-decision-chain} is the main modeling idea of the thesis.
First, a feasible timetable is selected. Second, feasible trip-to-trip
continuations are written as a graph. Third, a path cover of this graph gives
vehicle blocks. Fourth, the blocks are assigned to EV or ICE and checked for
charging and capacity feasibility. The final cost is computed only after the
complete schedule exists.

\section{Compatibility-Graph Formulation}
\label{sec:compatibility-graph-formulation}

\subsection{Trip Nodes and Time-Ordered Candidate Pairs}

For a fixed selected timetable \(\pi\), every selected trip becomes a graph
node. Only time-ordered trip pairs can become graph arcs:

\begin{equation}
P^{time}_\pi
=
\{(i,j)\in \pi\times\pi: et(i)<st(j)\}.
\label{eq:time-ordered-pairs}
\end{equation}

The compatibility arc set is a subset of these ordered pairs:

\begin{equation}
A_\pi\subseteq P^{time}_\pi.
\label{eq:arc-subset-time}
\end{equation}

The graph is therefore directed and acyclic under the time ordering. This
property is important because every path in the graph can be interpreted as a
vehicle block.

\subsection{In-Line Continuation Edges}

An in-line edge is possible when the end node of the first trip is the start
node of the second trip. The available stop duration is

\begin{equation}
\Delta_{ij}=st(j)-et(i).
\label{eq:inline-stop-duration}
\end{equation}

The in-line compatibility rule is

\begin{equation}
en(i)=sn(j),
\qquad
\delta^{min}_{en(i),\omega(et(i))}
\le
\Delta_{ij}
\le
\delta^{max}_{en(i),\omega(et(i))}.
\label{eq:inline}
\end{equation}

The time window is selected by the arrival time at the terminal, \(et(i)\).
This follows the problem statement: if a break crosses several time windows,
the stopping-time bounds are those of the first window, i.e., the arrival
window at the node \cite{minoa-problem}.

If \(\chi^{in}_{ij}=1\) denotes in-line compatibility, then

\begin{equation}
\chi^{in}_{ij}
=
\mathbf{1}
\left[
en(i)=sn(j)
\;\wedge\;
\delta^{min}_{en(i),\omega(et(i))}
\le \Delta_{ij}\le
\delta^{max}_{en(i),\omega(et(i))}
\right].
\label{eq:inline-indicator}
\end{equation}

\subsection{Depot-Bridge Continuation Edges}

If two trips cannot be connected directly at the same terminal, the vehicle can
return to the depot and leave again. This creates a depot-bridge edge. With
depot node \(dep\), the feasibility condition is

\begin{equation}
en(i)\neq sn(j),
\qquad
et(i)
+\tau^{PI}_{en(i),\omega(et(i))}
+\delta^{min}_{dep,\omega(et(i)+\tau^{PI}_{en(i),\omega(et(i))})}
+\tau^{PO}_{sn(j),\omega(st(j))}
\le
st(j).
\label{eq:depot}
\end{equation}

This inequality is the basic temporal test for a depot bridge. It does not by
itself describe all EV charging side conditions. If a vehicle stops at a
terminal before pull-in or after pull-out, that stop is allowed only for
charging and must also satisfy the corresponding charging, maximum stopping
time, and capacity rules \cite{minoa-problem}. In the implementation, a depot
bridge is written as pull-in, possible depot break, and pull-out. Terminal
charging adjacent to pull-in or pull-out is not inserted by the current
heuristic.

The corresponding indicator is

\begin{equation}
\begin{aligned}
\chi^{dep}_{ij}
=
\mathbf{1}\big[&
en(i)\neq sn(j)
\;\wedge\;
et(i)
+\tau^{PI}_{en(i),\omega(et(i))}\\
&+\delta^{min}_{dep,\omega(et(i)+\tau^{PI}_{en(i),\omega(et(i))})}
+\tau^{PO}_{sn(j),\omega(st(j))}
\le st(j)
\big].
\end{aligned}
\label{eq:depot-indicator}
\end{equation}

If the instance does not define a positive depot stopping time, then
\(\delta^{min}_{dep}=0\). The term remains in the model so that the condition is
correct for general input instances with a positive depot stop.

\subsection{Arc Set and Edge Attributes}

A directed arc is accepted if at least one of the two compatibility tests is
true:

\begin{equation}
(i,j)\in A_\pi
\Longleftrightarrow
(i,j)\in P^{time}_\pi
\;\wedge\;
\left(\chi^{in}_{ij}=1 \;\vee\; \chi^{dep}_{ij}=1\right).
\label{eq:arc-acceptance}
\end{equation}

For each accepted arc, the implementation stores a small vector of operational
attributes:

\begin{equation}
\xi_{ij}
=
\left(
t^{dead}_{ij},
t^{wait}_{ij},
t^{break}_{ij},
d^{drive}_{ij},
\text{type}_{ij}
\right).
\label{eq:arc-attribute-vector}
\end{equation}

These values are used by the greedy and weighted path-cover algorithms. The
reported objective is still computed later from the full vehicle schedule.

\begin{figure}[H]
\centering
  \includegraphics[width=\textwidth]{figures/fig30_feasible_trip_graph.pdf}
\caption{Compatibility-graph formulation.}
\label{fig:compatibility-example}
\end{figure}
\FloatBarrier

\Cref{fig:compatibility-example} gives the intuition behind
\Cref{eq:arc-acceptance}. The graph is not a drawing of the physical network.
It is a graph of possible vehicle continuations. Each selected trip is a node.
An arc is added when the same vehicle can serve the second trip after the first
one. In-line arcs use \Cref{eq:inline}, depot-bridge arcs use
\Cref{eq:depot}, and the final acceptance rule is given in
\Cref{eq:arc-acceptance}. Later, the variables \(x_{ij}\) and \(z_{ij}\)
select arcs from this graph.  The main message is that vehicle scheduling is
reduced to choosing compatible continuations, while all incompatible pairs are
removed before matching is solved.  This compatibility graph connects the
mathematical formulation to the later matching algorithm and should be read as
an operational graph, not as a geographic map.

\section{Decision Variables and State Variables}

\subsection{Timetable Selection Variables}

The selected timetable is denoted by \(\pi\). Its binary variable form is

\begin{equation}
s_i=
\begin{cases}
1, & i\in \pi,\\
0, & i\notin \pi.
\end{cases}
\label{eq:trip-selection-variable}
\end{equation}

For a direction \(d\), selected trips satisfy

\begin{equation}
\pi_d=\{i\in T_d:s_i=1\},
\qquad
\pi=\{i\in T:s_i=1\}.
\label{eq:selected-set-from-variable}
\end{equation}

\subsection{Coverage and Chain Variables}

The coverage variable links trips to vehicle blocks:

\begin{equation}
y_{ib}=
\begin{cases}
1, & \text{trip }i\text{ is covered by block }b,\\
0, & \text{otherwise.}
\end{cases}
\label{eq:coverage-variable}
\end{equation}

The direct continuation variable is

\begin{equation}
x_{ij}=
\begin{cases}
1, & \text{trip }j\text{ directly follows trip }i,\\
0, & \text{otherwise.}
\end{cases}
\label{eq:chain-variable}
\end{equation}

This variable is the algebraic version of a selected arc in
\Cref{fig:compatibility-example}. It enters the predecessor and successor limits
in \Cref{eq:successor-limit,eq:predecessor-limit}.

The matching model uses \(z_{ij}\) for the same predecessor-successor choice in
the bipartite graph:

\begin{equation}
z_{ij}=1
\Longleftrightarrow
(i^L,j^R)\text{ is selected in the matching}.
\label{eq:matching-variable}
\end{equation}

The variable \(z_{ij}\) is introduced separately because the implemented
path-cover algorithms solve a bipartite matching problem. It appears in the
unweighted matching objective \Cref{eq:max-cardinality-matching}, the matching
constraints \Cref{eq:matching-left,eq:matching-right,eq:matching-binary}, and
the weighted matching objective \Cref{eq:max-weight-matching}.

\subsection{Vehicle-Type Assignment Variables}

For each block \(b\in B\), the vehicle-type decision is represented by the
general variable

\begin{equation}
u_{bv}\in\{0,1\}
\qquad
\forall b\in B,\; v\in V.
\label{eq:fleet-variable-domain}
\end{equation}

Exactly one vehicle type is assigned:

\begin{equation}
\sum_{v\in V}u_{bv}=1
\qquad
\forall b\in B.
\label{eq:fleet-one-type}
\end{equation}

In the tested Senior instances, \(V\) contains one electric type and one ICE
type. For result reporting I therefore still count EV and ICE blocks, but the
mathematical notation does not rely on there being only two vehicle types.

\subsection{Battery, Charging, and Break Variables}

For each activity \(a\) in a block \(b\), \(r_{ba}^{pre}\) and
\(r_{ba}^{post}\) denote residual EV autonomy before and after the activity.
Charging durations are split by mode, and \(p_a\) denotes paid-break duration:

\begin{equation}
r_{ba}^{pre},r_{ba}^{post}\ge 0,\qquad
q_a^{slow},q_a^{fast}\ge 0,\qquad p_a\ge 0.
\label{eq:state-variable-domains}
\end{equation}

These variables are not solved in one large exact model in the implementation.
They are computed and checked during EV assignment and output construction.
They connect the software checks to the equations that follow. Residual autonomy
is propagated in \Cref{eq:battery-drive,eq:battery-charge,eq:battery-bounds}. Charging
duration appears in
\Cref{eq:charge-within-break,eq:charge-upper-slow,eq:charge-upper-fast,eq:model-min-charge}.
The paid-break variable \(p_a\) is used in the objective term
\Cref{eq:break-cost} and is defined more explicitly in
\Cref{eq:paid-break-inline,eq:paid-break-other}.

\section{Objective Function}

\subsection{Reported Vehicle-Scheduling Cost}

All reported costs follow the MINOA vehicle-scheduling objective. For a set of
vehicle blocks \(B\), the vehicle-scheduling cost is denoted by
\(C^{VS}\). The result tables report the same quantity as the cost reported by
the MINOA desktop validator for the final output.

\begin{equation}
C^{VS}
=
\sum_{b\in B}
\sum_{v\in V}
\left[
c_v^{fix}
+ c^{break} t^{break}_{b}
+ c_v^{pio}\left(t_b^{pi}+t_b^{po}\right)
+ c_v^{\mathrm{CO}_2}d(b)
\right]u_{bv}.
\label{eq:objective}
\end{equation}

Here \(c_v^{fix}\) is the fixed cost of vehicle type \(v\),
\(t_b^{break}\) is paid break time, \(t_b^{pi}\) is pull-in time, and
\(t_b^{po}\) is pull-out time. In the reported objective decomposition,
\(d(b)\) denotes the passenger-service driving duration to which the
instance-specific CO\(_2\) coefficient is applied:

\begin{equation}
d(b)
=
\sum_{i\in b}\tau_i^{trip}.
\label{eq:model-driving-time}
\end{equation}

Although the written MINOA problem description refers to \(d(b)\) as the total
driving time of block \(b\), the evaluation convention used for the reported
outputs applies the CO\(_2\) coefficient to passenger-service trip duration.
Pull-in and pull-out durations are accounted for separately through the
\(c_v^{pio}\) term. The formulation and the reported values follow this
convention, and the internal objective reconstruction agrees with the external
audit values.

This is different from EV autonomy, which is measured in kilometres and
decreases for both passenger trips and deadhead movements.
The CO\(_2\) coefficient \(c_v^{\mathrm{CO}_2}\) is zero for EV types and positive for ICE
types when specified in the input data \cite{minoa-problem,minoa-io}.

\subsection{Objective Decomposition}

For interpretation, I decompose the objective into four parts:

\begin{equation}
C^{VS}
=
C^{fix}+C^{break}+C^{pull}+C^{\mathrm{CO}_2}.
\label{eq:cost-decomposition}
\end{equation}

The fixed vehicle cost is

\begin{equation}
C^{fix}
=
\sum_{b\in B}\sum_{v\in V}c_v^{fix}u_{bv}.
\label{eq:fixed-cost}
\end{equation}

The paid-break cost is

\begin{equation}
C^{break}
=
c^{break}\sum_{b\in B}\sum_{\alpha\in A_b^{break}} p_\alpha.
\label{eq:break-cost}
\end{equation}

Pull-in and pull-out cost is

\begin{equation}
C^{pull}
=
\sum_{b\in B}
\sum_{v\in V}
c_v^{pio}
\left(t_b^{pi}+t_b^{po}\right)
u_{bv}.
\label{eq:pull-cost}
\end{equation}

The conventional-emission cost is

\begin{equation}
C^{\mathrm{CO}_2}
=
\sum_{b\in B}
\sum_{v\in V}
c_v^{\mathrm{CO}_2}d(b)u_{bv}.
\label{eq:co2-cost}
\end{equation}

This decomposition separates improvements from fewer vehicles, shorter breaks,
lower pull movement, and a different EV/ICE split.

\subsection{Internal Edge Scores}

The algorithms also use local edge scores, but these scores are not the
reported objective. For an arc \((i,j)\), a generic local score is

\begin{equation}
\hat c_{ij}
=
\alpha_1t^{dead}_{ij}
+\alpha_2t^{wait}_{ij}
+\alpha_3t^{break}_{ij}.
\label{eq:model-edge-score}
\end{equation}

This score helps rank feasible continuations. The final tables, however, use
\(C^{VS}\) from \Cref{eq:objective}.

\section{Core Feasibility Constraints}

\subsection{Timetable Feasibility Constraints}

For each direction \(d\), a directed headway graph \(H_d=(T_d,E_d)\) is built:

\begin{equation}
(i,j)\in E_d
\quad\Longleftrightarrow\quad
a(i)<a(j)
\;\wedge\;
a(j)-a(i)\le I^{max}_{ij}.
\label{eq:headway-edge}
\end{equation}

The achieved headway is

\begin{equation}
w_{ij}=a(j)-a(i),
\label{eq:headway}
\end{equation}

and every consecutive selected pair satisfies

\begin{equation}
w_{ij}\le I^{max}_{ij}.
\label{eq:headway-bound}
\end{equation}

The selected timetable of direction \(d\) is an initial-to-final path
\(P_h^d\) in this graph. The full selected timetable is

\begin{equation}
\pi=\bigcup_{d\in D}\pi_d,
\qquad
\pi_d=P_h^d .
\label{eq:selected-trips}
\end{equation}

For the sparse timetable construction, the dynamic recursion is

\begin{equation}
D_j=1+\min_{i:(i,j)\in E_d}D_i,
\label{eq:timetable-dp}
\end{equation}

with \(D_j=1\) for admissible initial trips.

\subsection{Timetable--Vehicle Schedule Linking}

Every selected trip is covered once:

\begin{equation}
\sum_{b\in B} y_{ib}=1
\qquad
\forall i\in \pi.
\label{eq:trip-covered-once}
\end{equation}

Unselected trips are not covered:

\begin{equation}
\sum_{b\in B} y_{ib}=0
\qquad
\forall i\in T\setminus \pi.
\label{eq:unselected-not-covered}
\end{equation}

In path-cover form, a selected trip can have at most one direct successor and
at most one direct predecessor:

\begin{align}
\sum_{j:(i,j)\in A_\pi} x_{ij} &\le 1
\qquad &&\forall i\in \pi, \label{eq:successor-limit}\\
\sum_{i:(i,j)\in A_\pi} x_{ij} &\le 1
\qquad &&\forall j\in \pi. \label{eq:predecessor-limit}
\end{align}

\subsection{Matching Constraints for Path Cover}

The path-cover model uses a bipartite graph with left trip copies \(\pi^L\)
and right trip copies \(\pi^R\). The unweighted matching objective is

\begin{equation}
\max
\sum_{(i,j)\in A_\pi} z_{ij}.
\label{eq:max-cardinality-matching}
\end{equation}

The matching constraints are

\begin{align}
\sum_{j:(i,j)\in A_\pi} z_{ij} &\le 1
\qquad &&\forall i\in \pi, \label{eq:matching-left}\\
\sum_{i:(i,j)\in A_\pi} z_{ij} &\le 1
\qquad &&\forall j\in \pi, \label{eq:matching-right}\\
z_{ij} &\in \{0,1\}
\qquad &&\forall (i,j)\in A_\pi. \label{eq:matching-binary}
\end{align}

If \(m=\sum z_{ij}\) continuations are selected, then the number of vehicle
blocks is

\begin{equation}
|B|=|\pi|-m .
\label{eq:path-cover-count}
\end{equation}

The weighted path-cover variant keeps the same constraints and uses

\begin{equation}
\max
\sum_{(i,j)\in A_\pi} w_{ij}z_{ij}.
\label{eq:max-weight-matching}
\end{equation}

\begin{figure}[H]
\centering
  \includegraphics[width=\textwidth]{figures/fig32_bipartite_matching_pathcover.pdf}
\caption{Bipartite matching transformation for path cover.}
\label{fig:bipartite-path-cover}
\end{figure}
\FloatBarrier

\Cref{fig:bipartite-path-cover} visualizes
\Cref{eq:max-cardinality-matching,eq:matching-left,eq:matching-right,eq:path-cover-count}.
The left copy of a trip can be matched to at most one successor, and the right
copy can be matched to at most one predecessor. After the matching is computed,
the selected arcs are followed to reconstruct vehicle blocks. A trip with no
matched predecessor starts a new block. A trip with no matched successor ends a
block. This is the reason why the number of blocks is \(|\pi|-|M_G|\).  The
figure therefore links a standard graph algorithm to the operational idea of
reducing the number of vehicles.  The two copies of the selected timetable do
not represent two physical networks; they represent predecessor and successor
positions in a vehicle block.

\subsection{Path-Cover Theorem for the Fixed Compatibility Graph}

The matching model above is exact for one fixed selected timetable and one
fixed compatibility graph. The following proposition states the graph result
used by the implementation.

\begin{proposition}[Minimum path cover by bipartite matching]
\label{prop:path-cover-matching}
Let \(G_\pi=(\pi,A_\pi)\) be the directed compatibility graph of a selected
timetable \(\pi\). Assume that all arcs respect time order, so \(G_\pi\) is
acyclic. Let \(\mathcal{B}_\pi=(\pi^L,\pi^R,E^B)\) be the bipartite graph with
one left and one right copy of every trip and with
\[
(i^L,j^R)\in E^B
\quad\Longleftrightarrow\quad
(i,j)\in A_\pi .
\]
If \(M^\star\) is a maximum-cardinality matching in \(\mathcal{B}_\pi\), then
the minimum number of vertex-disjoint directed paths covering all trips in
\(\pi\) is
\begin{equation}
K_\pi^\star
=
|\pi|-|M^\star|.
\label{eq:path-cover-theorem}
\end{equation}
\end{proposition}

\noindent\textit{Proof.}
Every selected matching edge \((i^L,j^R)\) can be interpreted as the
predecessor-successor relation \(i\rightarrow j\). Because a matching uses each
left copy and each right copy at most once, every trip has at most one selected
successor and at most one selected predecessor. Since all arcs point forward in
time, these selected arcs cannot form a directed cycle. They therefore form a
set of directed paths. Starting with \(|\pi|\) isolated trips, each selected
matching edge joins two paths and reduces the path count by one. A matching of
size \(|M|\) therefore gives \(|\pi|-|M|\) paths. Conversely, any path cover
with \(K\) paths contains exactly \(|\pi|-K\) predecessor-successor links, and
these links form a matching in \(\mathcal{B}_\pi\). Minimizing \(K\) is
therefore equivalent to maximizing \(|M|\), which proves
\Cref{eq:path-cover-theorem}.

\Cref{prop:path-cover-matching} is important because it explains why matching
is the exact subproblem inside the path-cover algorithms. Its scope is limited:
it minimizes the number of blocks only for the fixed graph \(G_\pi\). It does
not prove optimality of timetable selection, vehicle-type assignment, charging
insertion, capacity use, or the full MINOA objective.

\subsection{Vehicle-First Weighted Matching}

The weighted path-cover variant uses the same feasible edge set but prefers
operationally shorter continuations among maximum-cardinality matchings. For
an accepted arc \((i,j)\), the local transition cost is

\begin{equation}
\hat c_{ij}
=
t_{ij}^{dh}
+t_{ij}^{wait}
+\eta_{dep}t_{ij}^{dep},
\label{eq:weighted-local-cost}
\end{equation}

where \(t_{ij}^{dh}\) is deadhead time, \(t_{ij}^{wait}\) is non-depot waiting
time, \(t_{ij}^{dep}\) is depot waiting time, and \(0<\eta_{dep}<1\) gives
depot waiting a smaller local penalty. In the implementation
\(\eta_{dep}=0.1\). The matching edge weight is

\begin{equation}
w_{ij}=M^{veh}-\hat c_{ij},
\qquad M^{veh}=86400.
\label{eq:weighted-edge-main}
\end{equation}

The value \(M^{veh}\) is one day in seconds. In addition, the implementation calls
the matching routine with a maximum-cardinality requirement. Thus the
cardinality priority does not rely only on the numerical size of \(M^{veh}\). The
optimization is best read lexicographically:

\begin{equation}
\max |M|
\quad\text{first, then}\quad
\min \sum_{(i,j)\in M}\hat c_{ij}.
\label{eq:weighted-lexicographic}
\end{equation}

The local edge weight is a construction score. It is not the official MINOA
objective. After a full candidate schedule has been reconstructed, the
complete cost \(C^{VS}\) from \Cref{eq:objective} is computed internally. The
MINOA desktop validator is then used as an audit for the final reported
input-output pair.

\section{Electric Fleet, Charging, and Capacity Constraints}

\subsection{Electric Fleet Availability}

The number of blocks assigned to each electric vehicle type is bounded by the
available fleet:

\begin{equation}
\sum_{b\in B}u_{bv}\le N_v
\qquad
\forall v\in V_E.
\label{eq:ev-fleet-limit}
\end{equation}

ICE vehicles are treated as unbounded in the Senior data
\cite{minoa-problem,minoa-io}. In the tested instances, this means the ICE
count is obtained after EV assignment as

\begin{equation}
\sum_{b\in B}u_{b,ICE}
=
|B|-\sum_{b\in B}\sum_{v\in V_E}u_{bv}.
\label{eq:ice-fleet-count}
\end{equation}

\subsection{Battery Propagation}

For an EV block \(b\) assigned to vehicle type \(v\in V_E\), residual autonomy
is measured in kilometres. The vehicle leaves the depot fully charged:

\begin{equation}
r_{ba_1}^{pre}=A_v^{max}.
\label{eq:battery-initial}
\end{equation}

For a driving activity \(a\), which can be either a passenger trip or a
deadhead movement, autonomy decreases by travelled distance:

\begin{equation}
r_{ba}^{post}
=
r_{ba}^{pre}
-
\ell_a.
\label{eq:battery-drive}
\end{equation}

For a charging activity, autonomy increases by the charged amount:

\begin{equation}
r_{ba}^{post}
=
r_{ba}^{pre}
+
\Delta r_{ba}^{charge}.
\label{eq:battery-charge}
\end{equation}

The autonomy bounds are

\begin{equation}
0
\le
r_{ba}^{pre},r_{ba}^{post}
\le
A_v^{max}
\qquad
\forall b\in B,\;a\in A_b,\;v\in V_E.
\label{eq:battery-bounds}
\end{equation}

\subsection{Charging Feasibility}

Charging duration must fit into the available break:

\begin{equation}
q_a^{slow}+q_a^{fast}\le \tau_a
\qquad
\forall a\in A_b^{break}.
\label{eq:charge-within-break}
\end{equation}

The implementation uses minute-granular charging windows in the JSON output:

\begin{equation}
q_a^m=60k_a^m,
\qquad
k_a^m\in\mathbb{Z}_{\ge0},
\qquad
m\in\{\mathrm{slow},\mathrm{fast}\}.
\label{eq:minute-granularity}
\end{equation}

Slow charging cannot exceed the time needed to reach full autonomy:

\begin{equation}
q_a^{slow}
\le
T_v^{slow}
\frac{A_v^{max}-r_{ba}^{pre}}{A_v^{max}}.
\label{eq:charge-upper-slow}
\end{equation}

For fast charging, the full-charge time is \(\phi_vT_v^{slow}\):

\begin{equation}
q_a^{fast}
\le
\phi_vT_v^{slow}
\frac{A_v^{max}-r_{ba}^{pre}}{A_v^{max}}.
\label{eq:charge-upper-fast}
\end{equation}

The autonomy gained from charging is

\begin{equation}
\Delta r_{ba}^{charge}
=
A_v^{max}
\left(
\frac{q_a^{slow}}{T_v^{slow}}
+
\frac{q_a^{fast}}{\phi_vT_v^{slow}}
\right).
\label{eq:charge-gain}
\end{equation}

If charging is used, the minimum charging-time rule is

\begin{equation}
q_a^m=0
\quad\text{or}\quad
q_a^m\ge q_v^{min,m},
\qquad
m\in\{\mathrm{slow},\mathrm{fast}\}.
\label{eq:model-min-charge}
\end{equation}

At most one charging mode is active during one break window. The implementation
therefore writes either a slow-charging window, a fast-charging window, or a
non-charging parking window \cite{minoa-io}.

\subsection{Parking and Charging Capacity}

The output distinguishes the occupied resource from active charging. Let

\begin{equation}
o_a^{park},o_a^{slow},o_a^{fast},g_a\in\{0,1\}.
\label{eq:spot-occupation-vars}
\end{equation}

For each break window, exactly one resource is occupied:

\begin{equation}
o_a^{park}+o_a^{slow}+o_a^{fast}=1.
\label{eq:one-spot-occupied}
\end{equation}

Active charging is possible only while occupying a charging spot:

\begin{equation}
g_a\le o_a^{slow}+o_a^{fast}.
\label{eq:active-charging-spot}
\end{equation}

Let \(\mathcal{A}_n(t)\) be the set of break windows active at node \(n\) at
time \(t\). Capacity counts occupied slots, not only active charging:

\begin{align}
\sum_{a\in\mathcal{A}_n(t)}o_a^{park} &\le K_n^{park}, \label{eq:parking-capacity}\\
\sum_{a\in\mathcal{A}_n(t)}o_a^{slow} &\le K_n^{slow}, \label{eq:slow-capacity}\\
\sum_{a\in\mathcal{A}_n(t)}o_a^{fast} &\le K_n^{fast}. \label{eq:fast-capacity}
\end{align}

The JSON field \texttt{typeSpot} records which resource is occupied, while
\texttt{isCharging} records whether active charging occurs in that window
\cite{minoa-io}. A vehicle may therefore occupy a charging spot with
\texttt{isCharging=false} after charging has finished, but the spot still
counts against charging-slot capacity.

\begin{figure}[H]
\centering
  \includegraphics[width=\textwidth]{figures/fig35_capacity_ledger.png}
\caption{Capacity ledger at one terminal node.}
\label{fig:capacity-ledger}
\end{figure}
\FloatBarrier

\Cref{fig:capacity-ledger} gives the operational meaning of
\Cref{eq:parking-capacity,eq:slow-capacity,eq:fast-capacity}. The method does
not only check a vehicle at its arrival time. It checks interval overlap. For
example, a slow-charging interval uses both time and charger capacity until it
ends. If another vehicle tries to use the same charger during an overlapping
interval, the schedule is capacity infeasible.  The figure also explains why
the implementation uses a ledger: each inserted stop reserves a resource over a
time interval, and later stops must be checked against this reservation.  This
distinction is important for EV scheduling because a charging decision is
feasible only if the charger is available for the whole occupied interval.

\section{Paid-Break Accounting and Linearization}

\subsection{Direct Paid-Break Calculation}

For a stop activity \(a\), let \(\tau_a\) be total stop time,
\(\tau_a^{charge}=q_a^{slow}+q_a^{fast}\) be active charging time, \(n(a)\)
be the node, and \(w(a)=\omega(s_a)\) be the time window selected by the start
of the stop. For an in-line stop between two passenger trips, the paid-break
time is

\begin{equation}
p_a
=
\max\left\{
0,\;
\tau_a-\max\left(\tau_a^{charge},\delta^{min}_{n(a),w(a)}\right)
\right\}.
\label{eq:paid-break-inline}
\end{equation}

For stops adjacent to pull-in or pull-out movements, the official rule is
different. Such a terminal stop is permitted only as a charging stop; it does
not create positive paid-break time. The model description therefore uses

\begin{equation}
p_a
=
\begin{cases}
\begin{aligned}
\max\{0,\,
&\tau_a-\max(\tau_a^{charge},\\
&\delta^{min}_{n(a),w(a)})\},
\end{aligned}
&
\text{if \(a\) is an in-line stop},\\[1ex]
0,
&
\text{if \(a\) is adjacent to pull-in or pull-out}.
\end{cases}
\label{eq:paid-break-other}
\end{equation}

The implementation follows this distinction in the explanatory cost
decomposition: \texttt{costs.py} counts paid break only for breaks between two
passenger trips. The final output is then audited with the MINOA desktop
validator.

\begin{figure}[H]
\centering
\includegraphics[width=\textwidth]{figures/fig38_paid_break_accounting.png}
\caption{Paid-break accounting between two passenger trips.}
\label{fig:paid-break-accounting}
\end{figure}
\FloatBarrier

\Cref{fig:paid-break-accounting} shows how the waiting time between two trips is
split. Mandatory stopping time and active charging are not counted as paid
break. The remaining orange part is the cost-relevant break time.  This
interpretation is later used when result figures separate fixed vehicle cost
from smaller operational components.  Paid-break cost is therefore not the same
as total waiting time: mandatory stopping time and active charging are handled
first, and only the remaining waiting time contributes to \(C^{break}\).

\subsection{Big-M View of the Same Logic}

The direct calculation in \Cref{eq:paid-break-inline} can be written with the
auxiliary variable

\begin{equation}
m_a=\max\left(\tau_a^{charge},\delta^{min}_{n(a),w(a)}\right),
\label{eq:paid-break-max}
\end{equation}

and

\begin{equation}
p_a=\max\{0,\tau_a-m_a\}.
\label{eq:paid-break-positive}
\end{equation}

With a sufficiently large constant \(M^{break}\) and binary variables
\(\lambda_a,\eta_a\), one linearization is

\begin{align}
m_a &\ge \tau_a^{charge}, &
m_a &\ge \delta^{min}_{n(a),w(a)}, \notag\\
m_a &\le \tau_a^{charge} + M^{break}\lambda_a, &
m_a &\le \delta^{min}_{n(a),w(a)} + M^{break}(1-\lambda_a), \label{eq:paid-break-linearization-a}\\
p_a &\ge \tau_a-m_a, &
p_a &\ge 0, \notag\\
p_a &\le \tau_a-m_a+M^{break}(1-\eta_a), &
p_a &\le M^{break}\eta_a, \label{eq:paid-break-linearization-b}\\
\tau_a-m_a &\le M^{break}\eta_a, &
\tau_a-m_a &\ge -M^{break}(1-\eta_a). \notag
\end{align}

The implemented solver does not solve this as a separate MIP. The equations show
that the direct accounting is consistent with an exact linearized model once the
activity intervals are fixed.

\section{Timetable Variant Scoring}

\subsection{Sparse Headway Paths}

The timetable variant module selects an initial-to-final path in every
direction. A path \(P\) is first ranked by its number of selected trips. Among
paths with equal cardinality, a start-specific tie-breaking score is used:

\begin{equation}
P_h^d
=
\arg\min_{P\in\mathcal{P}^d}
\left(|P|,Q_h(P)\right).
\label{eq:timetable-follow-selection}
\end{equation}

The score is

\begin{equation}
Q_h(P)
=
\sum_{(i,j)\in P}g^h_{ij}
+g^h_{start}(i_0)
+g^h_{end}(i_q),
\label{eq:timetable-follow-score}
\end{equation}

where \(i_0\) and \(i_q\) are the first and last selected trips of the path.
Equivalently, the generic timetable-variant problem is

\begin{equation}
P_h^d
=
\arg\min_P
\left(
|P|,
\sum_{(i,j)\in P}g^h_{ij}
\right).
\label{eq:timetable-variant}
\end{equation}

\subsection{Examples of Variant Scores}

The implementation uses several forms of \(g^h_{ij}\):

\begin{equation}
g^h_{ij}\in
\left\{
-a(j),\;
a(j),\;
a(j)-a(i),\;
-(a(j)-a(i)),\;
\left|(a(j)-a(i))-\lambda_h I^{max}_{ij}\right|
\right\}.
\label{eq:timetable-score-examples}
\end{equation}

These scores generate different but feasible timetables. The purpose is not to
add a passenger penalty to the reported objective. The purpose is to create
several feasible graph structures for the multi-start path-cover method.

\section{Model Size and Complexity Discussion}

\subsection{Graph Size}

Let \(n=|\pi|\) be the number of selected trips. In the worst case, every
earlier trip can be tested against every later trip:

\begin{equation}
|P^{time}_\pi|
\le
\frac{n(n-1)}{2}.
\label{eq:time-pair-bound}
\end{equation}

The compatibility graph is smaller because only feasible in-line or
depot-bridge pairs become arcs:

\begin{equation}
|A_\pi|\le |P^{time}_\pi|.
\label{eq:arc-bound}
\end{equation}

This explains why the method first selects a sparse timetable. A smaller
\(\pi\) reduces the graph size and therefore also the matching problem.

\subsection{Matching Problem Size}

The bipartite graph contains two copies of the selected timetable:

\begin{equation}
|\pi^L\cup\pi^R|=2|\pi|.
\label{eq:bipartite-node-count}
\end{equation}

Its edge set has one edge for each compatibility arc:

\begin{equation}
|E^B|=|A_\pi|.
\label{eq:bipartite-edge-count}
\end{equation}

Thus, the core optimization problem is controlled by the number of selected
trips and feasible continuations, not by the full candidate pool \(T\). In
practice, this reduces the graph size before matching is solved and keeps the
path-cover subproblem focused on trips that are actually present in the
selected timetable.

\subsection{Why the Matheuristic is Practical}

The exact integrated problem would have to decide timetable, vehicle blocks,
fleet type, charging, battery states, and capacity at the same time. The
implemented method separates this into smaller linked subproblems:

\begin{equation}
\begin{aligned}
&\text{headway path}
+\text{compatibility graph}\\
&\quad
+\text{matching path cover}
+\text{EV feasibility simulation}.
\end{aligned}
\label{eq:method-decomposition}
\end{equation}

\Cref{eq:method-decomposition} explains why the later algorithm is a
matheuristic. The method does not prove global optimality, but it keeps the
main vehicle-count structure in the matching step. EV charging and capacity are
then checked on complete vehicle blocks, where their timing can be evaluated
directly.

\section{Exact Subproblems and Heuristic Coupling}
\label{sec:exact-subproblems-heuristic-coupling}

The formulation above contains two different levels of optimization.  One level
is exact after the timetable and compatibility graph have already been fixed.
The other level is heuristic because the algorithm must still choose which
timetable variant, edge score, fleet assignment, and charging insertion should
be retained.  Keeping these two levels separate is important for interpreting
the numerical results.  The matching step is solved as a well-defined graph
optimization problem, but the complete integrated MINOA problem is not solved to
global optimality.

\subsection{Fixed-Graph Path-Cover Layer}
\label{subsec:fixed-graph-path-cover-layer}

For one selected timetable \(\pi\), the compatibility graph
\(G_\pi=(\pi,A_\pi)\) is acyclic because every arc moves forward in time.  A
path cover of this graph is a set of directed paths that covers each selected
trip exactly once.  If \(z_{ij}=1\) means that trip \(j\) is served directly
after trip \(i\), the fixed-graph path-cover problem can be written as

\begin{align}
\max \quad
& \sum_{(i,j)\in A_\pi} w_{ij}z_{ij}
\label{eq:fixed-graph-path-cover}\\
\text{s.t.}\quad
& \sum_{j:(i,j)\in A_\pi} z_{ij}\le 1
&& \forall i\in\pi,
\label{eq:fixed-graph-successor}\\
& \sum_{i:(i,j)\in A_\pi} z_{ij}\le 1
&& \forall j\in\pi,
\label{eq:fixed-graph-predecessor}\\
& z_{ij}\in\{0,1\}
&& \forall (i,j)\in A_\pi .
\label{eq:fixed-graph-binary}
\end{align}

\Cref{eq:fixed-graph-successor} allows each trip to have at most one successor,
and \Cref{eq:fixed-graph-predecessor} allows each trip to have at most one
predecessor.  These two restrictions are exactly the rules needed to obtain
vehicle blocks from a matching.  Trips with no predecessor start a block, trips
with no successor end a block, and isolated trips become one-trip blocks.  The
number of blocks follows directly from the number of chosen continuations:

\begin{equation}
|B_\pi(z)|
=
|\pi|-\sum_{(i,j)\in A_\pi} z_{ij}.
\label{eq:block-number-from-continuations}
\end{equation}

\Cref{eq:block-number-from-continuations} is the main reason why matching is
attractive in this thesis.  Every accepted continuation reduces the vehicle
count by one, provided that the two trips can be linked in time and space.  In
the unweighted version, all feasible continuations are treated equally.  In the
weighted version, the score \(w_{ij}\) favors continuations that are also
promising with respect to waiting time, deadhead effort, EV feasibility, and
charging pressure.  This score is a construction score.  It is not the official
objective \(C^{VS}\), which is evaluated only after complete blocks and vehicle
types are known.

The fixed-graph layer can also be written as a bipartite matching problem.  The
left copy \(\pi^L\) represents possible predecessors and the right copy
\(\pi^R\) represents possible successors.  For every compatibility arc
\((i,j)\in A_\pi\), the bipartite graph contains the edge \((i^L,j^R)\).  A
matching then selects a set of trip-to-trip continuations without assigning two
successors to the same trip or two predecessors to the same trip.  This is the
computational form used by the implementation.

\subsection{Integrated Candidate Layer}
\label{subsec:integrated-candidate-layer}

The complete solution is not determined by one fixed graph.  It is selected
from a family of candidate graphs, edge scores, and fleet decisions.  Let
\(\mathcal{H}\) denote the set of timetable-variant generators and let
\(\Theta_h\) denote the bounded construction settings used for variant \(h\).
For each pair \((h,\theta)\), the implementation constructs a complete
candidate solution

\begin{equation}
S_{h,\theta}
=
\left(
\pi_h,\,
G_{\pi_h},\,
z_{h,\theta},\,
B_{h,\theta},\,
u_{h,\theta},\,
q_{h,\theta}
\right),
\label{eq:complete-candidate-tuple}
\end{equation}

where \(\pi_h\) is the selected timetable, \(G_{\pi_h}\) is the compatibility
graph, \(z_{h,\theta}\) is the chosen matching, \(B_{h,\theta}\) is the block
set, \(u_{h,\theta}\) assigns vehicle types, and \(q_{h,\theta}\) stores the
charging plan.  Only after all of these objects are available can the official
cost be computed:

\begin{equation}
C^{VS}(S_{h,\theta})
=
C^{fix}(S_{h,\theta})
+C^{break}(S_{h,\theta})
+C^{pull}(S_{h,\theta})
+C^{\mathrm{CO}_2}(S_{h,\theta}).
\label{eq:candidate-official-cost}
\end{equation}

The retained incumbent is therefore selected by

\begin{equation}
S^\star
\in
\arg\min
\left\{
C^{VS}(S_{h,\theta}) :
h\in\mathcal{H},\;
\theta\in\Theta_h,\;
S_{h,\theta}\in\mathcal{F}
\right\},
\label{eq:candidate-selection-rule}
\end{equation}

where \(\mathcal{F}\) is the set of internally feasible schedules.  A schedule
belongs to \(\mathcal{F}\) only if it satisfies timetable feasibility,
vehicle-block feasibility, EV autonomy, charging-duration rules, and parking or
charging capacity restrictions.  The MINOA desktop validator is then used as an
independent audit for the final output, not as part of
\Cref{eq:candidate-selection-rule}.

This separation also explains why charging-aware edge scoring is useful but
limited.  An edge score can estimate whether a continuation is likely to create
a good block, but battery feasibility is a property of the whole chronological
block.  A continuation that looks attractive locally may become unattractive
after several later trips are added.  Conversely, a slightly more expensive
continuation may create a block with better charging opportunities.  For this
reason, the final method combines local edge scores with full block-level
simulation before a candidate can enter \(\mathcal{F}\).

\subsection{Bound Interpretation}
\label{subsec:bound-interpretation-model}
\enlargethispage{2\baselineskip}

Every feasible candidate has an official cost and therefore gives an upper
bound on the unknown optimum:

\begin{equation}
C^\star \le C^{VS}(S_{h,\theta})
\qquad
\text{for every feasible } S_{h,\theta}.
\label{eq:upper-bound-from-candidate}
\end{equation}

Lower bounds have a different meaning.  A globally valid lower bound \(LB\)
must satisfy

\begin{equation}
LB \le C^\star \le UB,
\label{eq:global-bound-sandwich}
\end{equation}

where \(UB\) is the best feasible cost found.  Global optimality would be
certified only if \(LB=UB\).  In contrast, a selected-timetable diagnostic bound
can be useful for analysis but is conditional on the selected timetable
\(\pi\).  It therefore does not prove global optimality for the complete
integrated problem.  This distinction is used again in the computational
results, where the reported gaps are interpreted separately for global bounds
and selected-timetable diagnostic bounds.
